% =====================================================================
\chapter{Bài toán và mô tả tổng quan}
\label{ch:overview}
% =====================================================================

\section{Bối cảnh}
Trong quy trình gán nhãn hiện tại, reviewer kiểm tra annotation theo thứ tự job
hoặc frame. Cách làm này có ba hạn chế:

\begin{enumerate}
  \item \textbf{Lỗi phân bố không đều.} Phần lớn frame không có lỗi; lỗi tập
        trung ở cảnh đông đối tượng, đối tượng nhỏ, bị che, ban đêm, thời tiết
        xấu. Review tuần tự dồn phần lớn công sức vào frame sạch.
  \item \textbf{Thiếu bằng chứng tại chỗ.} Reviewer phải tự phát hiện lỗi thiếu
        box hoặc sai lớp bằng mắt; không có gợi ý nơi cần nhìn, guideline liên
        quan hay case đã phân xử.
  \item \textbf{Không đo được hiệu quả.} Không có reference độc lập nên không
        biết bao nhiêu lỗi còn sót, và không so sánh được cách làm mới với cách cũ.
\end{enumerate}

M13 giải quyết bài toán \emph{ưu tiên}: với ngân sách review hữu hạn, chọn
frame nào để xem trước sao cho bắt được nhiều lỗi thật nhất, đồng thời cho
reviewer bằng chứng đủ để ra quyết định nhanh và đúng. Quyết định cuối cùng
luôn thuộc về con người; trợ lý không tự xác nhận lỗi (R-02).

\section{Phát biểu bài toán}
\label{sec:problem}

\paragraph{Ký hiệu.}
\begin{itemize}
  \item \(F = \{f_1, \dots, f_N\}\): tập \(N\) frame trong phạm vi đánh giá (một snapshot).
  \item \(\mathcal{E}\): tập lỗi đã xác minh độc lập (reference) trên \(F\). Mỗi
        lỗi \(e \in \mathcal{E}\) thuộc đúng một frame \(\phi(e) \in F\) và một
        nhóm lỗi \(g(e) \in \{E1, E2, E3\}\).
  \item \(C_f\): tập candidate do các engine sinh ra trên frame \(f\); mỗi
        candidate \(c\) có đặc trưng \(\mathbf{x}_c\) và nhóm lỗi dự kiến \(g(c)\).
  \item \(s: F \to \mathbb{R}\): hàm điểm rủi ro; \(\pi\) là hoán vị của \(F\)
        theo \(s\) giảm dần, phá hoà theo khoá cố định.
  \item \(T_k(\pi)\): tập \(\lceil kN \rceil\) frame đầu của \(\pi\) với \(k \in (0,1]\).
\end{itemize}

\paragraph{Mục tiêu 1 — bắt lỗi sớm.}
Tìm \(s\) (thuộc họ hàm được chốt version trước khi đo) sao cho
\begin{equation}
  \mathrm{Recall@}k(\pi) \;=\;
  \frac{\bigl|\{e \in \mathcal{E} : \phi(e) \in T_k(\pi)\}\bigr|}{|\mathcal{E}|}
  \label{eq:recallk}
\end{equation}
đạt giá trị lớn tại \(k = 0{,}2\) và vượt có ý nghĩa thống kê so với các ranking
đối chứng (ngẫu nhiên, heuristic) trên tập held-out.

\paragraph{Mục tiêu 2 — giảm công sức với chất lượng tương đương.}
Gọi \(T_B\), \(T_A\) là tổng effort của quy trình baseline và assisted trên cùng
khối lượng dữ liệu; \(\rho_B\), \(\rho_A\) là tỉ lệ lỗi còn sót (residual) sau
mỗi quy trình. Mục tiêu:
\begin{equation}
  \text{giảm effort} \;=\; 1 - \frac{T_A}{T_B} \;\text{ lớn nhất có thể,}
  \quad\text{với ràng buộc}\quad \rho_A - \rho_B \le \delta .
  \label{eq:effort}
\end{equation}
\(T_A\) bắt buộc gồm thời gian xử lý cảnh báo sai và thời gian kiểm lại sau sửa.
Định nghĩa chi tiết, biên \(\delta\) và thiết kế thí nghiệm ở chương~\ref{ch:eval}.

\begin{notebox}[Vì sao đếm theo lỗi, không theo frame]
Một frame có thể chứa nhiều lỗi. Đếm theo frame (``frame có lỗi'') sẽ đánh giá
như nhau một frame có 1 lỗi và một frame có 6 lỗi. KPI của M13 đếm theo
\emph{lỗi đã xác minh}, đúng với đề bài. Recall theo frame có lỗi vẫn được báo
cáo như chỉ số phụ (mục~\ref{sec:kpi1}).
\end{notebox}

\section{Mục tiêu nghiệp vụ và chỉ số thành công}
\label{sec:kpis}

\begin{xltabular}{\linewidth}{@{}L{1.5cm}L{3.6cm}Y L{2.2cm}@{}}
\caption{Chỉ số thành công của M13}\label{tab:kpis}\\
\toprule
\textbf{Mã} & \textbf{Chỉ số} & \textbf{Định nghĩa ngắn} & \textbf{Ngưỡng đạt}\\
\midrule
\endfirsthead
\toprule \textbf{Mã} & \textbf{Chỉ số} & \textbf{Định nghĩa ngắn} & \textbf{Ngưỡng đạt}\\ \midrule
\endhead
\bottomrule
\endfoot
KPI-1 & Recall lỗi trong top 20\% frame & Số lỗi đã xác minh nằm trong 20\% frame đầu / tổng số lỗi đã xác minh của tập held-out. Đo \emph{khả năng xếp hạng}. Chỉ tính khi $|\mathcal{E}| \ge E_{min}$; dưới ngưỡng ghi ``không đủ mẫu'', không kết luận đạt/trượt. & \TBD[-K1]; $E_{min}$ \TBD[-K4]\\
KPI-2 & Mức giảm công sức review & \(1 - T_A/T_B\), effort gồm review, xử lý cảnh báo sai, phân xử, kiểm lại sau sửa. & \TBD[-K2]\\
KPI-2b & Recall phát hiện thực tế & Chỉ số phụ đo trong thí nghiệm KPI-2: số lỗi reference reviewer thực sự xác nhận / tổng số lỗi reference của phần dữ liệu đó, theo từng nhánh. & Báo cáo, không có ngưỡng\\
G-1 & Chất lượng tương đương & Non-inferiority một phía (\(\alpha = 0{,}025\)): cận trên khoảng tin cậy 95\% của \(\rho_A - \rho_B\) nhỏ hơn \(\delta\). & \(\delta\) \TBD[-K3]\\
G-2 & Gate: lỗi nghiêm trọng & Số issue mức nghiêm trọng đã xác nhận còn mở trên run cuối. & 0 (cấu hình pilot)\\
G-2b & Thí nghiệm: lỗi nghiêm trọng & Lỗi nghiêm trọng của reference còn lại (kể cả chưa có issue) ở assisted không nhiều hơn baseline. & Không nhiều hơn\\
G-5 & Thí nghiệm: box thừa & Box thừa (BR-06) trong annotation cuối của assisted không nhiều hơn baseline quá \(\delta_{fp}\). & \(\delta_{fp}\) \TBD[-K3]\\
G-3 & Đồng thuận reviewer & Số quyết định của reviewer trùng kết luận adjudicator / tổng quyết định được chấm, trên tập case chuẩn (theo H, PerformanceEvaluation). Quyết định ``Chưa chắc chắn'' tính là không trùng; mỗi reviewer tính riêng rồi báo cáo cả trung vị. Không thay thế accuracy của engine. & \(\ge 90\%\) (cấu hình pilot)\\
G-4 & Residual vận hành & Tỉ lệ frame còn lỗi ước lượng từ lát kiểm tra ngẫu nhiên (theo H), dùng khi không có reference toàn tập. Khác \(\rho\) của thí nghiệm. & \(\le 5\%\) (cấu hình pilot)\\
\end{xltabular}

\begin{decisionbox}
Ngưỡng G-2, G-3, G-4 là \emph{cấu hình khởi điểm pilot} đã chốt trong nguồn R
(Coverage 95\%, critical chưa đóng 0, rework verify 100\%, residual \(\le 5\%\),
agreement \(\ge 90\%\)). Đây chưa phải bằng chứng đạt. Ngưỡng KPI-1, KPI-2 được
đặt từ phép đo baseline trên tập hiệu chỉnh, \emph{trước} khi chạy held-out
(B-14).
\end{decisionbox}

\section{Bối cảnh hệ thống}
Hình~\ref{fig:context} thể hiện M13 trong môi trường: các tác nhân con người,
CVAT là hệ thống nguồn chỉ đọc, kho đối tượng lưu ảnh và evidence, và
Detector là mô hình độc lập chạy trên GPU worker.

\begin{figure}[H]
\centering
\begin{tikzpicture}[node distance=8mm and 14mm]
  \node[boxb, minimum width=46mm, minimum height=26mm, font=\sffamily\small\bfseries] (sys)
    {LabelX\\[2pt]\normalfont\scriptsize Module QC · M13 Reviewer Prioritization};
  % người
  \pic at (-6.2, 1.6) {actor=Reviewer};
  \pic at (-6.2,-1.2) {actor=Annotator};
  \pic at (6.2, 1.6) {actor={QA Lead}};
  \pic at (6.2,-1.2) {actor={QC Admin}};
  \pic at (0, 3.4) {actor={Product / Data Owner}};
  % hệ thống ngoài
  \node[ext, below left=18mm and 4mm of sys, text width=26mm] (cvat) {CVAT\\\scriptsize (nguồn annotation, editor)};
  \node[ext, below=18mm of sys, text width=24mm] (obj) {Object Storage\\\scriptsize (ảnh, evidence)};
  \node[ext, below right=18mm and 4mm of sys, text width=26mm] (det) {Detector baseline\\\scriptsize (GPU worker)};
  % luồng
  \draw[arr] (-5.7,1.9) -- node[lbl, above, sloped]{review, quyết định} (sys.west |- 0,0.6);
  \draw[arr] (sys.west |- 0,-0.4) -- node[lbl, below, sloped]{yêu cầu sửa, deep link} (-5.7,-0.9);
  \draw[arr] (5.7,1.9) -- node[lbl, above, sloped]{phân xử, reference} (sys.east |- 0,0.6);
  \draw[arr] (5.7,-0.9) -- node[lbl, below, sloped]{cấu hình, quyền} (sys.east |- 0,-0.4);
  \draw[arr] (0,2.6) -- node[lbl, right]{báo cáo hiệu quả} (sys.north);
  \draw[arr] (cvat.north) -- node[lbl, sloped, above]{đọc (chỉ đọc)} (sys.south west);
  \draw[arr, <->] (obj.north) -- (sys.south);
  \draw[arr, <->] (det.north) -- node[lbl, sloped, above]{inference} (sys.south east);
  \draw[arrd] (-6.2,-2.2) |- node[lbl, pos=0.75]{sửa annotation} (cvat.west);
\end{tikzpicture}
\caption{Sơ đồ ngữ cảnh hệ thống (context diagram)}
\label{fig:context}
\end{figure}

\section{Chức năng chính}
\begin{figure}[H]
\centering
\begin{tikzpicture}[node distance=5mm and 5mm]
  \node[boxg, text width=30mm] (m1) {\textbf{1. Thu thập}\\CVAT Adapter (chỉ đọc)\\Snapshot \& hash\\Kiểm drift};
  \node[boxg, text width=30mm, right=of m1] (m2) {\textbf{2. Phân tích nghi vấn}\\Schema/Taxonomy\\Geometry · Duplicate\\Detector + Matching};
  \node[boxg, text width=30mm, right=of m2] (m3) {\textbf{3. Hợp nhất \& xếp hạng}\\Candidate $\rightarrow$ Issue\\Risk score $s(f)$\\Ranking + audit slice};
  \node[boxg, text width=30mm, right=of m3] (m4) {\textbf{4. Review \& phân xử}\\Hàng đợi ưu tiên\\Workspace + evidence\\Escalation};
  \node[boxa, text width=30mm, below=9mm of m2] (s1) {\textbf{Hỗ trợ}\\Guideline lookup\\Rework + re-check\\Gate tối thiểu};
  \node[boxgr, text width=30mm, below=9mm of m4] (m5) {\textbf{5. Đo \& báo cáo}\\Reference \& Recall@k\\Effort log\\Báo cáo hiệu quả};
  \node[boxa, text width=30mm, below=9mm of m3] (s2) {\textbf{Nền tảng}\\RBAC · Lease\\Audit append-only\\Coverage ledger};
  \draw[arr] (m1)--(m2); \draw[arr] (m2)--(m3); \draw[arr] (m3)--(m4); \draw[arr] (m4)--(m5);
  \draw[arrd] (m4.south west) -- (s1.north east);
  \draw[arrd] (s1.west) -| (m1.south);
\end{tikzpicture}
\caption{Sơ đồ khối chức năng}
\label{fig:functions}
\end{figure}

\section{Lớp người dùng}
\label{sec:users}
\begin{xltabular}{\linewidth}{@{}L{3cm}Y L{3.6cm}@{}}
\caption{Lớp người dùng và quyền chính}\label{tab:users}\\
\toprule
\textbf{Vai trò} & \textbf{Mô tả và nhiệm vụ trong M13} & \textbf{Quyền chính}\\
\midrule
\endfirsthead
\toprule \textbf{Vai trò} & \textbf{Mô tả} & \textbf{Quyền chính}\\ \midrule
\endhead
\bottomrule
\endfoot
Reviewer & Người dùng chính. Làm việc theo hàng đợi ưu tiên, xem bằng chứng, ra quyết định, yêu cầu sửa, xác minh sau sửa. & Review; yêu cầu/xác minh rework. Không review annotation của mình.\\
Annotator & Sửa annotation trên CVAT theo yêu cầu rework; tham gia calibration. & Xem yêu cầu sửa của mình; thực hiện rework trên CVAT.\\
Quality Assurance Lead & Phân xử, quản lý nội dung guideline, chịu trách nhiệm reference và tập held-out; duyệt ngoại lệ khi được cấp quyền. & Adjudication; quản lý audit/calibration; duyệt waiver (khác người yêu cầu).\\
Quality Control Admin & Cấu hình kỹ thuật: engine, ngưỡng, policy version, mapping guideline, scope quyền. & Configuration (engine, ngưỡng, quyền). Không chạy phân tích (theo H: QA Lead chạy/xem).\\
Product / Data Owner & Chốt phạm vi, ngưỡng KPI, đọc báo cáo hiệu quả; quản lý dữ liệu BDD100K và artifact Detector. & Xem báo cáo; duyệt tiêu chí.\\
Super Admin & Quản trị hệ thống. Không là người duyệt mặc định; mọi ghi đè phải có lý do và gắn nhãn trong audit. & Ghi đè có lý do; vẫn chịu self-review và tách nhiệm vụ.\\
\end{xltabular}

\section{Môi trường vận hành}
\begin{itemize}
  \item \textbf{Máy chủ ứng dụng:} Linux, Django + DRF, PostgreSQL, Redis, Celery worker CPU.
  \item \textbf{GPU worker:} chạy Detector theo lô (batch inference). Phần cứng cụ thể \TBD[-02].
  \item \textbf{CVAT:} phiên bản đang triển khai, được pin sau khi kiểm API (\TBD[-01]).
  \item \textbf{Object Storage:} hệ thống hiện có (S3-compatible), lưu ảnh snapshot và evidence.
  \item \textbf{Client:} trình duyệt desktop hiện hành (Chrome/Edge/Firefox), độ phân giải tối thiểu 1440$\times$900.
\end{itemize}

\section{Ràng buộc thiết kế}
\label{sec:constraints}
Các quyết định sau đã chốt và ràng buộc thiết kế M13. Danh sách đầy đủ B-01…B-21
ở Phụ lục~\ref{app:decisions}.

\begin{xltabular}{\linewidth}{@{}L{1.2cm}L{4cm}Y@{}}
\caption{Ràng buộc thiết kế chính kế thừa từ quyết định đã chốt}\label{tab:constraints}\\
\toprule
\textbf{Mã} & \textbf{Chủ đề} & \textbf{Ràng buộc đối với M13}\\
\midrule
\endfirsthead
\toprule \textbf{Mã} & \textbf{Chủ đề} & \textbf{Ràng buộc}\\ \midrule
\endhead
\bottomrule
\endfoot
B-01 & Backend & Modular monolith + worker pool: Django + DRF, PostgreSQL, Celery + Redis, Object Storage. FastAPI chỉ khi đã có backend FastAPI được xác minh.\\
B-02 & Snapshot & Export từng job, chuẩn hoá, hash, kiểm drift trước khi khoá; lưu ảnh và metadata.\\
B-03 & Revision phát hành & Sau rework tạo snapshot + QC run cuối trên revision đã sửa; báo cáo trỏ run cuối.\\
B-04 & Coverage & Coverage theo từng engine với mẫu số áp dụng; không loại frame lỗi khỏi mẫu số.\\
B-06 & Geometry & Chỉ kiểm cấu trúc/bounds/diện tích; không suy biên vật thể khi thiếu tham chiếu.\\
B-07 & Mô hình độc lập & Một Detector baseline, freeze artifact/checksum, mapping lớp; Classifier chưa bật.\\
B-10 & Chạy lại & Candidate và Issue tách riêng; retry idempotent; dedup theo scope/object/nhóm lỗi.\\
B-12 & Quyền & Scope quyền, assignee tại snapshot, tách người yêu cầu/duyệt, audit append-only cùng transaction.\\
B-13 & Guideline & Tra cứu trực tiếp rule ID + version + section + case đã duyệt; không RAG.\\
B-17 & Trạng thái & Bảng transition/event/actor/guard dùng trạng thái của H.\\
B-18 & Ghi CVAT & Adapter chỉ đọc; sửa annotation qua deep link trên CVAT.\\
B-19 & Trạng thái công khai & Not checked/Partial/Failed không bao giờ thành đạt.\\
B-20, B-05 & Reference \& metric & Reference được duyệt và khoá trước khi đo; metric tính trên reference; không dùng model làm trọng tài.\\
B-21 & Matching & Matching là bước của Mô hình độc lập để sinh candidate; agreement không gọi là Precision/Recall.\\
\end{xltabular}

\section{Giả định và phụ thuộc}
\label{sec:assumptions}
\begin{xltabular}{\linewidth}{@{}L{1.3cm}Y L{3.4cm}@{}}
\caption{Giả định và phụ thuộc}\label{tab:assumptions}\\
\toprule
\textbf{Mã} & \textbf{Giả định / phụ thuộc} & \textbf{Nếu sai thì}\\
\midrule
\endfirsthead
\toprule \textbf{Mã} & \textbf{Giả định / phụ thuộc} & \textbf{Nếu sai thì}\\ \midrule
\endhead
\bottomrule
\endfoot
AS-01 & CVAT cho phép đọc job, frame, annotation, media qua API với token backend; xác định được thay đổi giữa hai lần đọc. & Không khoá được snapshot nhất quán; chặn toàn bộ luồng.\\
AS-02 & Có Detector baseline chạy được trên ảnh BDD100K với mapping đủ 10 lớp; artifact cố định được checksum. Pilot (v1.1): Faster R-CNN R-50-FPN 3x của model zoo BDD100K, huấn luyện trên tập \code{train}. & Engine Mô hình độc lập Not checked; không sinh được candidate E1/E2.\\
AS-03 & Detector không được huấn luyện trên ảnh của tập held-out (kiểm theo danh sách ảnh huấn luyện). Pilot (v1.1): danh sách ảnh huấn luyện là toàn bộ BDD100K \code{train}; ảnh thuộc \code{train} bị loại khỏi tập đánh giá. & Recall bị thổi phồng do rò rỉ dữ liệu; kết quả không hợp lệ.\\
AS-04 & QA Lead và ít nhất hai người xác minh độc lập sẵn sàng lập reference trước khi đo. Pilot (v1.1, CR-101): GT là nhãn gốc BDD100K; QA Lead duyệt và khoá reference, không cần người xác minh. & Không đo được KPI-1/KPI-2 (B-20).\\
AS-05 & Có ít nhất 4 reviewer tham gia thí nghiệm effort, trình độ tương đương. & Không cân bằng được hai nhánh; KPI-2 kém tin cậy.\\
AS-06 & Guideline gán nhãn BDD100K của dự án có rule ID và version. & Workspace chỉ hiện guideline dạng văn bản, không truy vết rule.\\
\end{xltabular}
